\begin{theorem}
For every $\varepsilon>0$ and every positive integer $k$, there is a natural
number $D_0$ such that the following holds for every real $D\ge D_0$ and
every $\theta\in(0,1)$. Let $\mathcal H$ be a finite $k$-uniform hypergraph
with every vertex degree at most $D$. If, for every real
$D'\ge\varepsilon D/3$ and every subhypergraph $H\subseteq\mathcal H$
with every vertex degree at most $D'$, every edge $e\in E(H)$ satisfies
\[
b^{\mathbb R}_{L(H),kD'}(e)\le\theta,
\]
where $b^{\mathbb R}$ is the real-scale local parameter, then
\[
\chi'(\mathcal H)\le(\theta+\varepsilon)kD.
\]
\end{theorem}
\begin{proof}
Suppose first that $0<\varepsilon\le1$, and choose the natural threshold
from \noderef{GeneralColoringIterationGreedyFinish}. Its main parameter
$D$ is real, but its local hypothesis only requires natural scales $n$.
For each natural $n\ge\varepsilon D/3$ and subhypergraph $H$ of maximum
degree at most $n$, apply the assumed real-scale hypothesis at the real
cast of $n$. By \noderef{LocalBParameterRealNatCast},
$b^{\mathbb R}_{L(H),kn}(e)=b_{L(H),kn}(e)$. Thus the local hypothesis of
the iteration lemma holds, and that lemma gives the conclusion.

If $\varepsilon>1$, choose $D_0=1$ and put $n=\lfloor D\rfloor_+$.
Then $n\ge1$, and \noderef{RealDegreeFloorBound} gives maximum degree at
most $n$. By \noderef{UniformHypergraphGreedyColoring}, there is a proper
edge-coloring with $kn$ colors. Since $n\le D$, $k>0$, and
$\theta+\varepsilon>1$, we have
$kn\le kD\le(\theta+\varepsilon)kD$.
This is the required real bound in \noderef{ChromaticIndexAtMostReal}.
The local hypothesis is not needed in this case.
\end{proof}
